\documentclass{report}
\usepackage{amsfonts, amsmath}
\newcommand\R{{\mathbb R}}
\newcommand\Q{{\mathbb Q}}
\newcommand\N{{\mathbb N}}
\pagestyle{empty}
\begin{document}
\large
\centerline{\Huge \bf Analisi Matematica I}
\vskip.2cm
\centerline{\Huge Correzione Primo Appello}

\vskip1cm
\begin{enumerate}
\item 
\[
-1-x<2x-1<1+x
\]
da cui $2>x>0$. 
\item L'insieme $A$ \`e  l'interno della circonferenza di raggio uno
centrata all'origine, $B$ invece la parte del piano che giace sopra la
bisettrice del primo e terzo quadrante. Il punto appartiene a $C$.
\item La funzione $f$ si chiama coseno iperbolico, su  qualunque libro
trovate il grafico. Due soluzioni visto che la retta $1+x$ sta sopra
il grafico di $f$ nel punto zero mentre cresce molto piu lentamente
all'infinito. Basta disegnare il grafico di $f(x)-1-x$ per averne la
controprova. 

\item La lunghezza che si deve minimizzare \`e data da
\[
\ell(x)=\sqrt{1+x^2}+\sqrt{1+(2-x)^2}.
\]
Calcolando la derivata e ponendola eguale a zero si ottiene
\[
x\sqrt{1+(2-x)^2}+(x-2)\sqrt{1+x^2}=0
\]
cio\`e
\[
x\sqrt{1+(2-x)^2}=-(x-2)\sqrt{1+x^2}
\]
prendendo il quadrato di ambo i membri e semplificando si ha $x=1$.

\item Poich\`e il numeratore \'e sempre minore di $n$ segue che il
limite \`e zero.

\item Si definisca $f$ nell'intervallo $[0,1]$ come $ax^3+bx^2+c x+d$.
Ora per avere la continuit\`a occorre $f(0)=0$ e $f(1)=1$, mentre per
avere la differenziabilit\`a occorre $f'(0)=0$ e $f'(1)=1$. da questo
segue: $d=0$, $c=0$, $a+b=1$, $3a+2b=1$. Cio\`e $a=-1$ e $b=2$.

\end{enumerate}
\end{document}
